\section{Introduction}

The progress of scientific discovery has always been intrinsically linked to our ability to measure physical phenomena with increasing precision and speed. Data Acquisition Systems (DAQ) serve as the vital bridge between the physical world and the digital domain of data analysis.

The performance of any experimental setup is fundamentally limited by the information collected, this process clearly begins with the detector. Detectors transform a physical interaction into a measurable electrical signal (typically voltage or current). However, modern experiments, ranging from high-energy particle physics to high-resolution medical imaging, generate a staggering volume of raw information. The challenge for modern DAQ systems is no longer just simple digitization, but the ability to handle high-bandwidth data streams while maintaining signal integrity and temporal synchronization \cite[pp. 668--669]{knollRadiationDetectionMeasurement2010a}.

In this context, Field Programmable Gate Arrays (FPGAs) have emerged as a transformative technology for scientific instrumentation. Unlike traditional General Purpose Processors (CPUs) or Digital Signal Processors (DSPs) that execute instructions sequentially, FPGAs provide a hardware-level parallel architecture. This allows for deterministic real-time processing, which is crucial for tasks such as "triggering"—the process of deciding in microseconds which data to keep and which to discard in high-frequency experiments [2]. Furthermore, the reconfigurability of FPGAs enables researchers to implement custom digital filters, pulse-height analyzers, and complex control loops directly into the hardware, significantly reducing the latency introduced by traditional software-based processing [3].